\noindent\textbf{Verbatim problem statement (from file).}
\emph{``Let $c>0$ and $h_c(n)$ be such that for any $n$ points in $\R^2$ such that there are $\ge cn^2$ lines each containing more than three points, there must be some line containing $h_c(n)$ many points. Estimate $h_c(n)$. Is it true that, for fixed $c>0$, we have $h_c(n)\to\infty$?''}

\medskip
\noindent\textbf{Definitions and conventions.}
Fix $c>0$.
For each $n\ge 1$, define $h_c(n)$ to be the smallest integer $m$ with the property:
\begin{quote}
For every point set $P\subset\R^2$ with $|P|=n$, if the number of distinct lines $\ell$ with $|\ell\cap P|\ge 4$ is at least $c n^2$, then there exists a line $\ell$ with $|\ell\cap P|\ge m$.
\end{quote}
The question is to estimate the growth of $h_c(n)$, and in particular whether for each fixed $c>0$ one necessarily has $h_c(n)\to\infty$ as $n\to\infty$.
